\section{Computational templates for reproducibility}\label{app:pseudocode}

\subsection{Single-point evaluation of \texorpdfstring{\(Z_m(y)\)}{Z\_m(y)}}

\begin{verbatim}
Input: m in N, y in C.

Initialize:
  Z0 = 1
  Z1 = y + 1
  Z2 = (y + 1)^2
  Z3 = (y + 1)^3

For k = 4,5,...,m:
  Zk = Z{k-1} + (2y+1) Z{k-2} - Z{k-3} - y(y+1) Z{k-4}

Output: Zm.
\end{verbatim}

\subsection{Certified isolation of the subdominant real branch point \texorpdfstring{\(\ySub\)}{ySub}}

Let \(c(y)=256y^3+411y^2+165y+32\). Since \(c\) has exactly one real root, any rational interval \((\alpha,\beta)\) with
\(c(\alpha)c(\beta)<0\) certifies \(\ySub\in(\alpha,\beta)\).

\begin{verbatim}
Input: rationals alpha < beta.
Compute c(alpha), c(beta).
If c(alpha)c(beta) < 0:
  (alpha, beta) is a certified enclosure of y_sub.
  Repeat bisection to refine to desired precision.
\end{verbatim}

\subsection{Dominant-modulus gap and zero-locus sampling}

For a fixed sample point \(y\in\CC\), compute the four roots of \(\Pi(\lambda,y)=0\),
sort them by decreasing modulus,
and record the dominant-modulus gap
\[
g(y)=|\lambda_{(1)}(y)|-|\lambda_{(2)}(y)|.
\]
Small values of \(g(y)\) approximate the dominant equimodular set predicted by the Beraha--Kahane--Weiss mechanism.
Separately, for a fixed \(m\), one may compute the roots of the polynomial \(Z_m(y)\) and compare them with the sampled small-gap region.

\begin{verbatim}
Input: grid points y in C, system size m.

For each sample point y:
  Compute the four roots lambda_i of Pi(lambda,y)=0.
  Sort them so |lambda_(1)| >= |lambda_(2)| >= |lambda_(3)| >= |lambda_(4)|.
  Record g(y) = |lambda_(1)| - |lambda_(2)|.

Compute the polynomial Z_m(y) from the recurrence and find its zeros.
Plot:
  - points where g(y) is small;
  - zeros of Z_m(y);
  - exact markers y = 0, y = 1, y = y_sub, y = -1.
\end{verbatim}

